% Options for packages loaded elsewhere
% Options for packages loaded elsewhere
\PassOptionsToPackage{unicode}{hyperref}
\PassOptionsToPackage{hyphens}{url}
\PassOptionsToPackage{dvipsnames,svgnames,x11names}{xcolor}
%
\documentclass[
  english,
  11pt,
  letterpaper,
]{article}
\usepackage{xcolor}
\usepackage[margin=1in]{geometry}
\usepackage{amsmath,amssymb}
\setcounter{secnumdepth}{-\maxdimen} % remove section numbering
\usepackage{iftex}
\ifPDFTeX
  \usepackage[T1]{fontenc}
  \usepackage[utf8]{inputenc}
  \usepackage{textcomp} % provide euro and other symbols
\else % if luatex or xetex
  \usepackage{unicode-math} % this also loads fontspec
  \defaultfontfeatures{Scale=MatchLowercase}
  \defaultfontfeatures[\rmfamily]{Ligatures=TeX,Scale=1}
\fi
\usepackage{lmodern}
\ifPDFTeX\else
  % xetex/luatex font selection
\fi
% Use upquote if available, for straight quotes in verbatim environments
\IfFileExists{upquote.sty}{\usepackage{upquote}}{}
\IfFileExists{microtype.sty}{% use microtype if available
  \usepackage[]{microtype}
  \UseMicrotypeSet[protrusion]{basicmath} % disable protrusion for tt fonts
}{}
\makeatletter
\@ifundefined{KOMAClassName}{% if non-KOMA class
  \IfFileExists{parskip.sty}{%
    \usepackage{parskip}
  }{% else
    \setlength{\parindent}{0pt}
    \setlength{\parskip}{6pt plus 2pt minus 1pt}}
}{% if KOMA class
  \KOMAoptions{parskip=half}}
\makeatother
% Make \paragraph and \subparagraph free-standing
\makeatletter
\ifx\paragraph\undefined\else
  \let\oldparagraph\paragraph
  \renewcommand{\paragraph}{
    \@ifstar
      \xxxParagraphStar
      \xxxParagraphNoStar
  }
  \newcommand{\xxxParagraphStar}[1]{\oldparagraph*{#1}\mbox{}}
  \newcommand{\xxxParagraphNoStar}[1]{\oldparagraph{#1}\mbox{}}
\fi
\ifx\subparagraph\undefined\else
  \let\oldsubparagraph\subparagraph
  \renewcommand{\subparagraph}{
    \@ifstar
      \xxxSubParagraphStar
      \xxxSubParagraphNoStar
  }
  \newcommand{\xxxSubParagraphStar}[1]{\oldsubparagraph*{#1}\mbox{}}
  \newcommand{\xxxSubParagraphNoStar}[1]{\oldsubparagraph{#1}\mbox{}}
\fi
\makeatother


\usepackage{longtable,booktabs,array}
\usepackage{calc} % for calculating minipage widths
% Correct order of tables after \paragraph or \subparagraph
\usepackage{etoolbox}
\makeatletter
\patchcmd\longtable{\par}{\if@noskipsec\mbox{}\fi\par}{}{}
\makeatother
% Allow footnotes in longtable head/foot
\IfFileExists{footnotehyper.sty}{\usepackage{footnotehyper}}{\usepackage{footnote}}
\makesavenoteenv{longtable}
\usepackage{graphicx}
\makeatletter
\newsavebox\pandoc@box
\newcommand*\pandocbounded[1]{% scales image to fit in text height/width
  \sbox\pandoc@box{#1}%
  \Gscale@div\@tempa{\textheight}{\dimexpr\ht\pandoc@box+\dp\pandoc@box\relax}%
  \Gscale@div\@tempb{\linewidth}{\wd\pandoc@box}%
  \ifdim\@tempb\p@<\@tempa\p@\let\@tempa\@tempb\fi% select the smaller of both
  \ifdim\@tempa\p@<\p@\scalebox{\@tempa}{\usebox\pandoc@box}%
  \else\usebox{\pandoc@box}%
  \fi%
}
% Set default figure placement to htbp
\def\fps@figure{htbp}
\makeatother



\ifLuaTeX
\usepackage[bidi=basic]{babel}
\else
\usepackage[bidi=default]{babel}
\fi
% get rid of language-specific shorthands (see #6817):
\let\LanguageShortHands\languageshorthands
\def\languageshorthands#1{}
\ifLuaTeX
  \usepackage[english]{selnolig} % disable illegal ligatures
\fi


\setlength{\emergencystretch}{3em} % prevent overfull lines

\providecommand{\tightlist}{%
  \setlength{\itemsep}{0pt}\setlength{\parskip}{0pt}}



 


\usepackage{needspace}
\makeatletter
\@ifpackageloaded{tcolorbox}{}{\usepackage[skins,breakable]{tcolorbox}}
\@ifpackageloaded{fontawesome5}{}{\usepackage{fontawesome5}}
\definecolor{quarto-callout-color}{HTML}{909090}
\definecolor{quarto-callout-note-color}{HTML}{0758E5}
\definecolor{quarto-callout-important-color}{HTML}{CC1914}
\definecolor{quarto-callout-warning-color}{HTML}{EB9113}
\definecolor{quarto-callout-tip-color}{HTML}{00A047}
\definecolor{quarto-callout-caution-color}{HTML}{FC5300}
\definecolor{quarto-callout-color-frame}{HTML}{acacac}
\definecolor{quarto-callout-note-color-frame}{HTML}{4582ec}
\definecolor{quarto-callout-important-color-frame}{HTML}{d9534f}
\definecolor{quarto-callout-warning-color-frame}{HTML}{f0ad4e}
\definecolor{quarto-callout-tip-color-frame}{HTML}{02b875}
\definecolor{quarto-callout-caution-color-frame}{HTML}{fd7e14}
\makeatother
\makeatletter
\@ifpackageloaded{caption}{}{\usepackage{caption}}
\AtBeginDocument{%
\ifdefined\contentsname
  \renewcommand*\contentsname{Table of contents}
\else
  \newcommand\contentsname{Table of contents}
\fi
\ifdefined\listfigurename
  \renewcommand*\listfigurename{List of Figures}
\else
  \newcommand\listfigurename{List of Figures}
\fi
\ifdefined\listtablename
  \renewcommand*\listtablename{List of Tables}
\else
  \newcommand\listtablename{List of Tables}
\fi
\ifdefined\figurename
  \renewcommand*\figurename{Figure}
\else
  \newcommand\figurename{Figure}
\fi
\ifdefined\tablename
  \renewcommand*\tablename{Table}
\else
  \newcommand\tablename{Table}
\fi
}
\@ifpackageloaded{float}{}{\usepackage{float}}
\floatstyle{ruled}
\@ifundefined{c@chapter}{\newfloat{codelisting}{h}{lop}}{\newfloat{codelisting}{h}{lop}[chapter]}
\floatname{codelisting}{Listing}
\newcommand*\listoflistings{\listof{codelisting}{List of Listings}}
\makeatother
\makeatletter
\makeatother
\makeatletter
\@ifpackageloaded{caption}{}{\usepackage{caption}}
\@ifpackageloaded{subcaption}{}{\usepackage{subcaption}}
\makeatother
\usepackage{bookmark}
\IfFileExists{xurl.sty}{\usepackage{xurl}}{} % add URL line breaks if available
\urlstyle{same}
\hypersetup{
  pdftitle={Repetition, Similarity, Relevance},
  pdfauthor={Jordan Gunn},
  pdflang={en},
  pdfkeywords={retrieved-context theory, repetition, semantic
similarity, narrative memory, research proposal},
  colorlinks=true,
  linkcolor={darkgray},
  filecolor={Maroon},
  citecolor={Blue},
  urlcolor={teal},
  pdfcreator={LaTeX via pandoc}}


\title{Repetition, Similarity, Relevance}
\usepackage{etoolbox}
\makeatletter
\providecommand{\subtitle}[1]{% add subtitle to \maketitle
  \apptocmd{\@title}{\par {\large #1 \par}}{}{}
}
\makeatother
\subtitle{Toward an Integrated Retrieved Context Account of Text
Comprehension and Free Recall}
\author{Jordan Gunn}
\date{February 2024}
\begin{document}
\maketitle
\begin{abstract}
I propose experiments to specify and evaluate competing explanations of
text comprehension and memory in terms of retrieved context theory
(RCT). RCT explains recall of word sequences by asserting that 1)
studied items are associated with a gradually changing, recency-weighted
representation of temporal context, 2) the contexts retrieved from
studying (or recalling) items update the current state of context, and
3) the current state of context is the proximate cue for each item
recall. RCT has been successful in accounting for a wide range of free
and serial recall data, explaining key benchmark phenomena, predicting
the effects of experimental manipulations, and generalizing to
naturalistic research conditions. Comprehension, the process of
constructing an integrated mental representation of the meaning of a
text, is known to engage many of the same cognitive processes as the
free and serial recall tasks, and preliminary evidence suggests that
text memory data exhibits many of the same patterns observed in free
recall. However, an integrated account of text comprehension and memory
in terms of RCT has remained elusive. Here, I argue that the main
obstacle to an integrated account is underspecification of how
similarity and feature overlap across study materials interact with
temporal context to structure the associations learned during study and
applied to guide recall. The proposed research will narrow this gap by
constructing parallel free recall datasets using stories and word lists
as study materials and manipulating relations between study events via
item repetition, item similarity, and item co-relevance. Through formal
comparison of candidate RCT-based models, the proposed research will
clarify key constraints for an integrated account of text comprehension
and memory to guide future research and theory development.
\end{abstract}


\begin{tcolorbox}[enhanced jigsaw, bottomrule=.15mm, breakable, toprule=.15mm, colframe=quarto-callout-note-color-frame, arc=.35mm, colback=white, rightrule=.15mm, opacityback=0, left=2mm, leftrule=.75mm]

\vspace{-3mm}\textbf{About this version}\vspace{3mm}

This is an earlier research proposal, preserved in February 2024 and
published here in October 2026. It records the original plan to connect
repetition, semantic reminding, and narrative memory. Subsequent
repetition-memory work challenged the neighbour-contiguity evidence used
as a starting point here, sharpening the question of when reinstatement
changes learning. The original proposal text is retained below,
including its provisional claims and citations.

This sketch is separate from the
\href{https://githubpsyche.github.io/vanderbilt-thesis-proposal-2024/}{full
Vanderbilt thesis proposal submitted in September 2024}, which is
available on its own site with separate chapter pages and the original
165-page PDF.

\end{tcolorbox}

\section{Introduction}\label{introduction}

Models of text comprehension and of serial memory share a common thread:
both concern how studying sequences of items come to be associated with
one another

\Needspace{8\baselineskip}

\subsection{Premise}\label{premise}

According to Retrieved Context Theory,

\begin{enumerate}
\def\labelenumi{\arabic{enumi}.}
\tightlist
\item
  Memory associates each study event in a list learning experiment with
  a state of temporal context, a representation that reflects your
  recent experiences.
\item
  When you study an item in a list-learning experiment, it reminds you
  of past experiences with that item. You retrieve its pre-existing
  contextual associations and integrate that into the current temporal
  context.
\end{enumerate}

These two points have interesting implications in a list with item
repetitions such that participants study some items at two positions,
\(i\) and \(j\). When you see the item at position \(j\), you reinstate
the temporal context from when you saw it at position \(i\). Because of
(1), a side effect of this reminder is that context will also reinstate
neighbors of the item's first presentation (\(i+1\), \(i+2\)). This
information will persist in context for a few study events, so the
item's first presentation neighbors get associated with the item's
second presentation neighbors. Even though the items at \(i+1\) and
\(i+2\) were never restudied!

Lohnas \& Kahana (2014) identified this implication and showed that CMR
predicts this neighbor contiguity effect and that the effect shows up in
behavioral data. When people later recall an item in \{\(i+1\),
\(i+2\)\}, they transition to recalling items in \{\(j+1\), \(j+2\)\} at
rates much higher than chance. Even if they never actually recall the
item repeated at positions \(i\) and \(j\).

This ties into semantic or narrative memory because this analysis helps
detect \emph{any} reminding effects in a free recall experiment. Lynn's
results suggest that when you see an item for the second time in a list
learning experiment, you tend to remember the first time you saw it in a
way that shapes the course of recall. We can ask if this happens when
you see \textbf{two distinct but highly related ideas} in a study
sequence.

If I see the word KING at position \(i\) and then the word QUEEN at
position \(j\) in a study list, do I get reminded of my experience of
the word KING in a way that shapes how I process the list going forward
and later remember it?

If I read a sentence about how dangerous wild animals are at one point
in a narrative and then a bear occurs as a character later, do I get
reminded of the ``wild animals'' sentence in a way that shapes how I
process the story going forward and later remember it?

I am sure a positive answer to the latter question is already common
knowledge in the narrative comprehension literature -- they call the
phenomenon
\href{https://en.wikipedia.org/wiki/Text_inferencing\#:~:text=Backward\%20inferences\%20require\%20the\%20reader\%20to\%20bridge\%20the\%20current\%20text\%20idea\%20to\%20one\%20that\%20occurred\%20earlier\%20in\%20the\%20text.}{bridging
inference} and make a big deal about it. I'd go as far as to say that
most model competition in the narrative comprehension literature
concerns when and how readers perform bridging inferences to comprehend
a text. For example, an ongoing debate in the narrative comprehension
literature distinguishes between strategy-based and memory-based
accounts of bridging inference (O'Brien \& Cook, 2015). Strategy-based
models assume that readers actively search the long-term memory
representation of the text for information relevant to the contents of
working memory. In contrast, memory-based models assume that inactive
information related to the current input is reactivated via a passive
``resonance mechanism''.

By generalizing Lohnas \& Kahana's (2014) neighbor contiguity analysis
to measure reminding effects driven by semantic similarity instead of
just item repetitions, we can obtain a way to relate the problem of
bridging inference in the narrative comprehension literature to parallel
issues surrounding semantic organization in free recall.

The models examined in Morton \& Polyn (2016) don't quite explore this
question because they don't allow semantic similarity to influence the
encoding process. Instead, the models all suppose that semantic
similarity between items only directly affects the recall process. They
thus don't predict a ``similar items'' neighbor contiguity effect that a
semantic resonance mechanism would enact.

I think with the right design, though, we can show that such an effect
does exist and that a straightforward modification to CMR can account
for it. Furthermore, I think we can use the same approach to show that
this same model can account for a range of bridging inference phenomena
in narrative comprehension. If successful, it would draw a direct line
from this core idea in narrative comprehension --- bridging inference
--- to a core idea in the list-learning literature --- context
reinstatement. And hopefully set the stage for a unified account of both
processes.

\subsection{Experiment 1: Demonstrating the
Principle}\label{experiment-1-demonstrating-the-principle}

To investigate whether similarity-driven contextual reinstatement
organizes the list learning process, we can start by replicating part of
Lynn's 2016 experiment.

Participants were presented with control lists and mixed lists. Each
list contained 40 unique positions. In the control lists, 40 unique
items were presented just once. In the mixed lists, 6 pairs of repeated
items were mixed among 28 once presented items. Across lists, there were
an equal number of item pairs repeated at spacings of lag \(\in\) \{0,
1, 2, \ldots, 8\}, where lag is defined as the number of intervening
items between an item's repetitions. 35 subjects saw lists of each type
12 times over 3 sessions.

To analyze the data, Lynn measured the percent of transitions between
items that follow the presentation of the same repeated item. If an item
was presented at positions \(i\) and \(j\), and if one of the items was
recalled from either \(S_j = \{j + 1, j + 2\}\) or
\(S_i = \{i + 1, i + 2\}\), they then counted the percentage of times
that a transition was then made to one of the items in the complementary
sets \(S_i\) or \(S_j\), respectively.

To measure the percent of transitions expected by chance, over 100
reshuffled datasets, repeated items in control lists were matched to the
same serial positions as the repeated items in mixed lists. Then the
percentage of transitions was counted.

We can replicate these results to measure similarity-driven contextual
reinstatement and add a condition similar to our mixed list condition.
In mixed lists, six pairs of repeated items were mixed among 28 once
presented items. In our additional condition, we can use six pairs of
items with high semantic similarity to one another but low semantic
similarity to other list items. We can enforce the remaining 28 items to
have comparatively low semantic similarity to all list items. Then to
analyze the data, we can measure the percent of transitions between
items that follow the presentation of items in the same pairing.

If we observe that this percent of transitions is significantly higher
than chance, we can conclude that similarity-driven contextual
reinstatement influences list learning and the course of free recall.

If we don't, but under the same conditions can replicate Lynn's results
regarding the mixed list condition, we can surmise that
similarity-driven contextual reinstatement does not influence list
learning.

It may be possible to do a similar demonstration with existing datasets,
but I figure a design like this would provide a clearer test of the
semantic resonance hypothesis.

\subsubsection{Modeling}\label{modeling}

Let's suppose do find a resonance effect. With data from Experiment 1
and by integrating a similarity-driven contextual reinstatement
mechanism into CMR, we can also test the hypothesis through model
comparison. We can demonstrate that the mechanism contributes unique
predictions about the outcome of our neighbor contiguity analysis in our
novel experimental condition. And through likelihood-based model
fitting, we can show that these differences either improve or worsen the
model's capacity to account for relevant datasets overall.

During list learning, the mechanism would alter how the model calculates
contextual input \(c^{IN}_i\). In the classic specification of CMR,
\(c^{IN}_i = M^{FC}f_i\), where \(M^{FC}\) is a matrix that learns
associations between orthogonalized item representations and contextual
states. In a modified model, we might maintain a parallel memory
structure \(M^{sem}\) that learns associations between distributed
semantic representations of items and contextual states. This way,
contextual reinstatement upon processing an item will be graded based on
the similarity of the item's semantic representation to semantic
representations of other items in memory.

I'm unsure, but these associations may prove easier to implement within
an instance-based architecture than a traditional linear associator
model. A nonlinear activation scaling mechanism could control the extent
to which the model reinstates similar items' contextual representations
upon item processing.

\subsection{Experiment 2: Generalization to Arbitrary Word
Lists}\label{experiment-2-generalization-to-arbitrary-word-lists}

A natural follow-up to our results in Experiment 1 is generalizing the
experiment to arbitrary word lists. Using benchmark datasets from Morton
\& Polyn (2016), we can perform a model comparison to evaluate whether
the hypothesized mechanism can generalize to improve CMR's account of
memory for word lists with less structured similarity relations between
studied words.

These results would help clarify the scope of results in Experiment 1,
ruling out or confirming that

\begin{enumerate}
\def\labelenumi{\arabic{enumi}.}
\tightlist
\item
  the design of our original experiment influenced subjects to process
  paired high-similarity words like item repetitions, but that subjects
  do not do so naturally.
\item
  similarity-driven contextual reinstatement can influence free recall
  even without enforced extremes in the semantic similarity between
  items that I propose we enact in Experiment 1
\end{enumerate}

\subsection{Experiment 3/4: Generalization to Written
Narratives}\label{experiment-34-generalization-to-written-narratives}

Finally, we can attempt to generalize our results to memory for written
narratives.

Two approaches are possible. On the one hand, we can try to create
narratives whose structures mirror those of the control, mixed, and
similarity-pairing lists from our Experiment 1. Each narrative can
contain a connected sequence of 40 ``idea units'' or subjects/predicates
in single words or phrases. In our ``control'' narratives, we'll enforce
idea units to be unique and have low similarity to one another. In our
``mixed'' narratives, 6 pairs of repeated idea units can be mixed among
28 once presented idea units. And in our similarity-pairing condition,
we can use 6 pairs of idea units with high semantic similarity to one
another but a low semantic similarity to other idea units. The remaining
28 idea units will be enforced to have comparatively low semantic
similarity to all idea units.

This design might enable a clear demonstration of a similarity-driven
contextual reinstatement mechanism in narrative recall but seems to
require constructing very unnatural stories since a typical narrative
includes some item repetition. So alternatively, we might use naturally
structured narratives identified from previous research into narrative
comprehension. Then we might apply the likelihood-based model comparison
technique from Morton \& Polyn (2016) to evaluate how including a
similarity-driven contextual reinstatement mechanism in CMR influences
the model's ability to account for narrative recall datasets.

At the same time, we'd also be evaluating CMR itself as an account of
narrative recall. We may include implementations of other comprehension
models in this analysis.

\subsection{Discussion}\label{discussion}

So what do we get from all this trouble? Ideally, a measure of
theoretical unification.

Similarity-driven contextual reinstatement is a powerful mechanism
wielded in models of narrative comprehension to account for readers'
ability to organize otherwise seemingly disjointed text ideas into a
coherent discourse representation. It may also be an important way
semantic knowledge shapes new episodic memories in general. These
proposed experiments suggest a way we might convincingly connect these
ideas under a single familiar theoretical framework -- retrieved context
theory.

Our extension of the neighbor contiguity analysis from Lohnas \& Kahana
(2016) gives us a way to cleanly demonstrate the behavioral predictions
suggested by the mechanism under controlled conditions and independently
of formal theoretical commitments. It also helps us focus on a specific
mechanism/theme instead of a general memory account as we go forward.

Similarly, the likelihood-based model comparison technique from Morton
\& Polyn (2016) gives us a way to evaluate the mechanism's effectiveness
as part of a complete account of memory search under more general
experimental conditions.

Using connected narratives rather than structured word lists to evaluate
our model allows us to connect our findings, methods, and ideas to
debates in the narrative comprehension literature and vice versa.

Through all this, we can validate our extension of CMR as an account not
just of memory search, but of story understanding. Our toolbox of free
recall analyses can begin to provide novel theoretical constraints for
the study of narrative comprehension. And in turn, a vast literature of
ideas and techniques exploring how rich semantic associations form can
inform the broader study of semantic and episodic memory.




\end{document}
